\documentclass[10pt,aspectratio=169]{beamer}

\usepackage{cmap}
\usepackage[T2A]{fontenc}
\usepackage[utf8]{inputenc}
\usepackage[russian]{babel}
\usepackage{amsmath,amssymb,mathtools}
\usepackage{booktabs}
\usepackage{graphicx}
\usepackage{array}
\usepackage{adjustbox}
\usepackage{hyperref}

\usetheme{Madrid}
\usecolortheme{seahorse}
\setbeamertemplate{navigation symbols}{}
\setbeamertemplate{blocks}[rounded][shadow=false]
\setbeamertemplate{title page}[default][colsep=-4bp,rounded=true,shadow=false]
\setbeamertemplate{footline}{%
  \leavevmode\hbox{%
    \begin{beamercolorbox}[wd=.80\paperwidth,ht=2.8ex,dp=1.2ex,leftskip=1em]{author in head/foot}%
      \usebeamerfont{author in head/foot}\insertshorttitle
    \end{beamercolorbox}%
    \begin{beamercolorbox}[wd=.20\paperwidth,ht=2.8ex,dp=1.2ex,rightskip=1em plus1fil]{date in head/foot}%
      \hfill\insertframenumber{} / \inserttotalframenumber
    \end{beamercolorbox}}%
}
\setbeamertemplate{caption}[numbered]
\hypersetup{unicode=true,colorlinks=true,linkcolor=blue,urlcolor=blue}

\graphicspath{{assets/labs_4_5/}}

\newcommand{\R}{\mathbb{R}}
\newcommand{\grad}{\nabla}
\newcommand{\norm}[1]{\left\lVert #1 \right\rVert}
\newcommand{\inner}[2]{\left\langle #1,#2 \right\rangle}
\newcommand{\eps}{\varepsilon}
\newcommand{\smallnote}[1]{\vspace{0.35em}{\footnotesize\color{gray!75!black}#1}}
\newcommand{\plotone}[2]{%
    \begin{center}
        \includegraphics[width=\linewidth,height=0.72\textheight,keepaspectratio]{#1}
    \end{center}
    \vspace{-0.4em}
    {\footnotesize #2}
}
\newcommand{\plottwo}[4]{%
    \begin{columns}[T,onlytextwidth]
        \column{0.49\textwidth}
        \includegraphics[width=\linewidth,height=0.60\textheight,keepaspectratio]{#1}
        \column{0.49\textwidth}
        \includegraphics[width=\linewidth,height=0.60\textheight,keepaspectratio]{#2}
    \end{columns}
    \vspace{0.2em}
    {\footnotesize \textbf{#3.} #4}
}

\title[ЛР4 и ЛР5]{Методы оптимизации: лабораторные работы 4 и 5}
\subtitle{Методы Ньютона, квазиньютоновские методы и data-driven задачи регрессии}
\author[А. Новиков, А. Душкин]{Алексей Новиков \and Александр Душкин}
\institute{Университет ИТМО · Курс «Методы оптимизации»}
\date{2026}

\begin{document}

\begin{frame}
    \titlepage
\end{frame}

\begin{frame}{Материалы и назначение презентации}
\begin{columns}[T,onlytextwidth]
\column{0.52\textwidth}
\textbf{Использованные материалы:}
\begin{itemize}
    \item условия ЛР4 и ЛР5;
    \item отчеты \texttt{Report\_demo4.tex}, \texttt{Report\_demo5.tex};
    \item код \texttt{functions.py}, \texttt{run.py}, \texttt{config.yaml};
    \item результаты из \texttt{outputs/lab4} и \texttt{outputs/lab5};
    \item конспект последних тем: CG, Newton, BFGS, GN, LM, regression.
\end{itemize}
\column{0.44\textwidth}
\textbf{Фокус защиты:}
\begin{itemize}
    \item постановка задачи и критерии остановки;
    \item какие методы реализованы и зачем;
    \item что показывают графики и таблицы;
    \item почему часть запусков формально неуспешна;
    \item выводы по методам и применимости.
\end{itemize}
\end{columns}
\end{frame}

\begin{frame}{Общие критерии сравнения}
\begin{block}{Критерий остановки}
\[
    \norm{\grad f(x_k)} \le 10^{-8}.
\]
\end{block}

\begin{columns}[T,onlytextwidth]
\column{0.48\textwidth}
\textbf{Метрики в ЛР4}
\begin{itemize}
    \item число итераций;
    \item число вызовов $f$;
    \item число вызовов $\grad f$;
    \item число вызовов $\grad^2 f$;
    \item причина остановки.
\end{itemize}
\column{0.48\textwidth}
\textbf{Метрики в ЛР5}
\begin{itemize}
    \item loss и empirical risk;
    \item число эпох или итераций;
    \item число вычислений градиента;
    \item время работы;
    \item коэффициенты модели.
\end{itemize}
\end{columns}

\smallnote{Успешный запуск означает выполнение строгого критерия остановки, а не просто хорошее значение функции.}
\end{frame}

\section{Лабораторная 4}

\begin{frame}{ЛР4: постановка и экспериментальный протокол}
\begin{columns}[T,onlytextwidth]
\column{0.5\textwidth}
\textbf{Задачи из условия}
\begin{itemize}
    \item линейный CG на SPD-квадратиках;
    \item нелинейные CG: Fletcher--Reeves и Polak--Ribiere;
    \item NewtonCholesky и NewtonDirectionChoice;
    \item trust-region Powell Dog Leg;
    \item DFP, BFGS, L-BFGS;
    \item сравнение с SciPy Newton-CG.
\end{itemize}
\column{0.47\textwidth}
\textbf{Эксперимент}
\begin{itemize}
    \item $n\in\{2,10,50,100\}$;
    \item $\kappa\in\{1,10,100,1000\}$;
    \item по два seed для каждой пары $(n,\kappa)$;
    \item 2D-траектории: quadratic, Rosenbrock, Himmelblau, Ackley;
    \item всего 552 запуска.
\end{itemize}
\end{columns}
\end{frame}

\begin{frame}{Квадратичная модель и обусловленность}
\[
    f(x)=\frac12\inner{Hx}{x}+\inner{c}{x}+b,
    \qquad \grad f(x)=Hx+c,\qquad \grad^2 f(x)=H.
\]

\begin{block}{Если $H\succ0$}
\[
    x^*=-H^{-1}c,\qquad
    f(x)=f(x^*)+\frac12\inner{H(x-x^*)}{x-x^*}.
\]
\end{block}

\begin{block}{Число обусловленности}
\[
    \kappa(H)=\frac{\lambda_{\max}}{\lambda_{\min}}.
\]
\end{block}

\smallnote{Рост $\kappa$ делает линии уровня вытянутыми. Методы без явной информации о кривизне начинают идти зигзагами или требуют больше линейных поисков.}
\end{frame}

\begin{frame}{CG и nonlinear CG}
\begin{columns}[T,onlytextwidth]
\column{0.49\textwidth}
\textbf{Линейный CG для SPD-квадратики}
\[
    f(x)=\frac12 x^TAx-b^Tx,\qquad A\succ0.
\]
\[
    p_0=-g_0,\qquad
    \alpha_k=\frac{g_k^Tg_k}{p_k^TAp_k},
\]
\[
    x_{k+1}=x_k+\alpha_kp_k,
    \qquad
    p_{k+1}=-g_{k+1}+\beta_kp_k.
\]
\column{0.49\textwidth}
\textbf{Нелинейные варианты}
\[
    \beta_k^{FR}=
    \frac{g_{k+1}^Tg_{k+1}}{g_k^Tg_k},
\]
\[
    \beta_k^{PR}=
    \frac{g_{k+1}^T(g_{k+1}-g_k)}{g_k^Tg_k}.
\]
\[
    \beta_k=\max(0,\beta_k^{PR})\quad\text{в PR+}.
\]
\end{columns}

\smallnote{Гарантия ``не более $n$ шагов'' относится к SPD-квадратике и точной арифметике. На нелинейных функциях это уже эвристика направления.}
\end{frame}

\begin{frame}{Newton, LM и trust-region}
\[
    m_k(p)=f(x_k)+g_k^Tp+\frac12p^TH_kp.
\]

\begin{columns}[T,onlytextwidth]
\column{0.48\textwidth}
\textbf{Newton}
\[
    H_kp_k=-g_k.
\]
Если $H_k\succ0$, то
\[
    g_k^Tp_k=-g_k^TH_k^{-1}g_k<0.
\]

\textbf{LM / сдвиг Гессиана}
\[
    (H_k+\lambda_k I)p_k=-g_k.
\]
\column{0.48\textwidth}
\textbf{Trust-region}
\[
    \min_{\norm{p}\le\Delta_k} m_k(p).
\]
\[
    \rho_k=
    \frac{f(x_k)-f(x_k+p_k)}
         {m_k(0)-m_k(p_k)}.
\]
Powell Dog Leg использует ломаную: шаг Коши $\to$ ньютоновский шаг.
\end{columns}
\end{frame}

\begin{frame}{Квазиньютоновские методы}
\[
    s_k=x_{k+1}-x_k,\qquad y_k=g_{k+1}-g_k,\qquad G_k\approx H_k^{-1}.
\]

\begin{block}{Секущее условие}
\[
    G_{k+1}y_k=s_k,\qquad y_k^Ts_k>0.
\]
\end{block}

\begin{columns}[T,onlytextwidth]
\column{0.49\textwidth}
\textbf{BFGS}
\[
G_{k+1}=(I-\rho sy^T)G_k(I-\rho ys^T)+\rho ss^T,
\quad \rho=\frac1{y^Ts}.
\]
\column{0.49\textwidth}
\textbf{DFP}
\[
G_{k+1}=G_k+\frac{ss^T}{s^Ty}
-\frac{G_kyy^TG_k}{y^TG_ky}.
\]
\end{columns}

\smallnote{L-BFGS не хранит полную матрицу $G_k$, а использует последние $m$ пар $(s_i,y_i)$, поэтому память $O(mn)$ вместо $O(n^2)$.}
\end{frame}

\begin{frame}{ЛР4: что реализовано в коде}
\begin{columns}[T,onlytextwidth]
\column{0.48\textwidth}
\textbf{Функции}
\begin{itemize}
    \item \texttt{GeneratedQuadratic};
    \item \texttt{Quadratic2DVisualization};
    \item \texttt{Lab4Rosenbrock};
    \item \texttt{Lab4Himmelblau};
    \item \texttt{Lab4Ackley}.
\end{itemize}
\column{0.49\textwidth}
\textbf{Особенности эксперимента}
\begin{itemize}
    \item history траекторий включает стартовую точку $x_0$;
    \item на Ackley и Himmelblau есть запуски около седловых областей;
    \item для NewtonDirectionChoice сохранялась диагностика выбора направления;
    \item для L-BFGS исследовались $m=3,5,10$.
\end{itemize}
\end{columns}
\end{frame}

\begin{frame}{ЛР4: средние метрики на GeneratedQuadratic}
\scriptsize
\begin{adjustbox}{max width=\textwidth}
\begin{tabular}{lrrrrr}
\toprule
Метод & Успешных & Итерации & $N_f$ & $N_g$ & $N_H$\\
\midrule
NewtonCholesky & 32 & 1.00 & 3.00 & 3.00 & 1.00\\
NewtonDirectionChoice & 32 & 1.00 & 3.00 & 3.00 & 1.00\\
PowellDogLeg & 32 & 4.94 & 5.94 & 5.94 & 4.94\\
QuadraticConjugateGradient & 16 & 70.25 & 71.25 & 71.25 & 1.00\\
BFGS & 17 & 31.09 & 240.12 & 65.06 & 0.00\\
DFP & 16 & 27.97 & 200.28 & 58.31 & 0.00\\
LBFGS & 40 & 68.99 & 261.59 & 140.66 & 0.00\\
ScipyNewtonCG & 25 & 10.69 & 28.66 & 27.22 & 10.88\\
FletcherReeves & 8 & 239.91 & 2997.56 & 498.41 & 0.00\\
PolakRibiere & 8 & 148.78 & 1727.84 & 307.00 & 0.00\\
\bottomrule
\end{tabular}
\end{adjustbox}

\smallnote{L-BFGS считается по трем значениям памяти $m$, поэтому успешных запусков может быть больше, чем число отдельных квадратичных задач.}
\end{frame}

\begin{frame}{ЛР4: влияние размерности и обусловленности}
\plottwo{lab4_n_iter_vs_n_k_10.png}
        {lab4_n_iter_vs_k_n_10.png}
        {Число итераций}
        {При росте $n$ и $\kappa$ методы без точного Гессиана сильнее накапливают цену. Ньютоновские методы почти не увеличивают число итераций, но их шаг дороже из-за линейной алгебры.}
\end{frame}

\begin{frame}{ЛР4: вычислительная цена}
\plottwo{lab4_n_grad_calls_vs_n_k_10.png}
        {lab4_n_hessian_calls_vs_n_k_10.png}
        {Вызовы производных}
        {Квазиньютоновские и CG-методы не вызывают Гессиан, но могут тратить больше градиентов. Ньютоновские методы экономят итерации, но платят вычислением и факторизацией Гессиана.}
\end{frame}

\begin{frame}{ЛР4: сравнение методов и память L-BFGS}
\plottwo{lab4_comparison_n_iter.png}
        {lab4_lbfgs_memory_n_iter.png}
        {Итоговое сравнение}
        {SciPy Newton-CG служит устойчивой библиотечной точкой сравнения. У L-BFGS параметр $m$ задает объем сохраняемой кривизны: больше памяти не всегда означает заметно меньше итераций.}
\end{frame}

\begin{frame}{ЛР4: диагностика NewtonDirectionChoice}
\scriptsize
\begin{adjustbox}{max width=\textwidth}
\begin{tabular}{lrrrrr}
\toprule
Функция & Запусков & newton & modified\_newton & steepest\_descent & max $\lambda$\\
\midrule
GeneratedQuadratic & 32 & 32 & 0 & 0 & 0\\
Lab4Ackley & 3 & 9 & 2 & 0 & 19.6791\\
Lab4Himmelblau & 3 & 11 & 3 & 0 & 42\\
Lab4Rosenbrock & 3 & 62 & 0 & 0 & 0\\
Quadratic2DVisualization & 5 & 5 & 0 & 0 & 0\\
\bottomrule
\end{tabular}
\end{adjustbox}

\begin{itemize}
    \item На SPD-квадратиках всегда достаточно обычного ньютоновского направления.
    \item На Ackley и Himmelblau появляются сдвиги $H+\lambda I$: локальная кривизна не всегда положительно определена.
    \item Переход к steepest descent не потребовался: сдвиг Гессиана оказался достаточным.
\end{itemize}
\end{frame}

\begin{frame}{ЛР4: траектории на квадратичной функции}
\plotone{Quadratic2DVisualization__seed_42_x0_-8.0_-6.0_x_star_1.0_-1.0__trajectory_panels.png}
{На квадратичной функции ньютоновское направление быстро попадает в минимум. CG- и квазиньютоновские методы постепенно восстанавливают геометрию задачи.}
\end{frame}

\begin{frame}{ЛР4: траектории на Rosenbrock}
\plotone{Lab4Rosenbrock__x0_-1.2_1.0__trajectory_panels.png}
{У Rosenbrock узкая изогнутая долина. Траектории могут делать дуги и корректировки; это нормальное поведение для методов, использующих локальную информацию о кривизне.}
\end{frame}

\begin{frame}{ЛР4: траектории на Himmelblau}
\plotone{Lab4Himmelblau__x0_3.35_0.08__trajectory_panels.png}
{Старт выбран около седловой области, но не ровно в седле. Методы могут уходить к разным областям притяжения, поэтому важно видеть стартовую точку и первые шаги.}
\end{frame}

\begin{frame}{ЛР4: траектории на Ackley}
\plotone{Lab4Ackley__x0_0.7_0.65__trajectory_panels.png}
{Ackley содержит локальные колебания. Траектория может визуально ``крутиться'', если локальная модель несколько раз меняет направление; это не ошибка само по себе.}
\end{frame}

\begin{frame}{ЛР4: выводы для защиты}
\begin{itemize}
    \item На SPD-квадратиках методы Ньютона ожидаемо выигрывают по итерациям: точный Гессиан полностью задает форму функции.
    \item Powell Dog Leg делает больше шагов, зато контролирует доверие к локальной квадратичной модели.
    \item BFGS и L-BFGS являются компромиссом: вторые производные не нужны, но информация о кривизне накапливается постепенно.
    \item Нелинейные CG чувствительны к линейному поиску и потере сопряженности.
    \item Формально успешный запуск требует $\norm{\grad f(x_k)}\le 10^{-8}$; остановка по \texttt{max\_iter} или \texttt{step\_tol} может дать хорошую точку, но не считается строгой сходимостью.
\end{itemize}
\end{frame}

\section{Лабораторная 5}

\begin{frame}{ЛР5: постановка и экспериментальный протокол}
\begin{columns}[T,onlytextwidth]
\column{0.5\textwidth}
\textbf{Задачи из условия}
\begin{itemize}
    \item сгенерировать near-linear и nonlinear данные;
    \item построить полиномиальные модели степеней 1--5;
    \item сравнить analytic solution, SGD, mini-batch, GN, LM;
    \item изучить batch size;
    \item сравнить L1, L2 и Elastic Net.
\end{itemize}
\column{0.47\textwidth}
\textbf{Эксперимент}
\begin{itemize}
    \item 120 точек в каждом датасете;
    \item шум: $0.04$ для near-linear и $0.09$ для nonlinear;
    \item $\lambda\in\{10^{-3},10^{-2},10^{-1},1\}$;
    \item batch size: $1,4,8,16,32,64,120$;
    \item всего 264 запуска.
\end{itemize}
\end{columns}
\end{frame}

\begin{frame}{Regression как задача оптимизации}
\begin{block}{Data-driven постановка}
\[
    \{(x_i,y_i)\}_{i=1}^{N},\qquad
    f(x,w)\approx y.
\]
\end{block}

\begin{block}{Empirical risk для регрессии}
\[
    Q(w)=\frac1N\sum_{i=1}^{N}(f(x_i,w)-y_i)^2.
\]
\end{block}

\begin{block}{Полиномиальная модель}
\[
    f(x,w)=\sum_{j=0}^{d} w_j z^j,
    \qquad
    z=\frac{x-\bar x}{s_x}.
\]
\end{block}

\smallnote{Нормировка $x$ уменьшает плохую обусловленность степеней $x^j$.}
\end{frame}

\begin{frame}{Лосс, риск и регуляризация}
\begin{columns}[T,onlytextwidth]
\column{0.5\textwidth}
\textbf{Loss на одном объекте}
\[
    \ell_i(w)=(f(x_i,w)-y_i)^2.
\]

\textbf{Средний риск}
\[
    Q(w)=\frac1N\sum_i \ell_i(w).
\]
\column{0.47\textwidth}
\textbf{Регуляризованная цель}
\[
    F(w)=Q(w)+\lambda_1\norm{w}_{1}+\lambda_2\norm{w}_2^2.
\]
\[
    |w_j|\approx\sqrt{w_j^2+\delta^2}
\]
для сглаженной L1.
\end{columns}

\smallnote{В экспериментах свободный коэффициент $w_0$ не регуляризируется: он отвечает за вертикальный сдвиг, а не за сложность формы.}
\end{frame}

\begin{frame}{Аналитическое решение и стохастические методы}
\begin{columns}[T,onlytextwidth]
\column{0.49\textwidth}
\textbf{1D linear regression}
\[
    f(x)=wx+b.
\]
\[
    w=
    \frac{\overline{xy}-\bar x\bar y}
         {\overline{x^2}-(\bar x)^2},
    \qquad
    b=\bar y-w\bar x.
\]

\column{0.49\textwidth}
\textbf{SGD / mini-batch}
\[
    w_{t+1}=w_t-\alpha_t \grad \ell_{i_t}(w_t),
\]
\[
    w_{t+1}=w_t-\alpha_t\frac1{|B_t|}
    \sum_{i\in B_t}\grad\ell_i(w_t).
\]
\end{columns}
\end{frame}

\begin{frame}{Gauss--Newton и Levenberg--Marquardt}
\[
    F(w)=\frac12\norm{r(w)}^2,\qquad
    r_i(w)=f(x_i,w)-y_i.
\]
\[
    \grad F(w)=J(w)^Tr(w),
\]
\[
    \grad^2 F(w)=J(w)^TJ(w)+\sum_{i=1}^{N}r_i(w)\grad^2 r_i(w).
\]

\begin{columns}[T,onlytextwidth]
\column{0.49\textwidth}
\textbf{Gauss--Newton}
\[
    J_k^TJ_kp_k=-J_k^Tr_k.
\]
\column{0.49\textwidth}
\textbf{Levenberg--Marquardt}
\[
    (J_k^TJ_k+\lambda_k I)p_k=-J_k^Tr_k.
\]
\end{columns}

\smallnote{Если residual мал или модель почти линейна по параметрам, $J^TJ$ хорошо приближает точный Гессиан.}
\end{frame}

\begin{frame}{ЛР5: что реализовано в коде}
\begin{columns}[T,onlytextwidth]
\column{0.48\textwidth}
\textbf{Данные и модель}
\begin{itemize}
    \item \texttt{near\_linear\_true};
    \item \texttt{nonlinear\_true};
    \item генерация шума;
    \item нормировка признака;
    \item матрица дизайна для степени $d$;
    \item \texttt{PolynomialRegressionObjective}.
\end{itemize}
\column{0.49\textwidth}
\textbf{Методы}
\begin{itemize}
    \item \texttt{AnalyticalLinearRegression1D};
    \item \texttt{StochasticGradientDescent};
    \item \texttt{MiniBatchGradientDescent};
    \item \texttt{GaussNewton};
    \item \texttt{LevenbergMarquardt};
    \item сохранение histories, coefficients, datasets, summary.
\end{itemize}
\end{columns}
\end{frame}

\begin{frame}{ЛР5: сравнение основных методов}
\scriptsize
\begin{adjustbox}{max width=\textwidth}
\begin{tabular}{lrrrrr}
\toprule
Метод & Успешных & Лучшее $f$ & Средн. итераций & $N_f$ & $N_g$\\
\midrule
AnalyticalLinearRegression1D & 2 & 0.0394 & 1.00 & 1.00 & 1.00\\
GaussNewton & 10 & 0.0385 & 1.00 & 2.00 & 2.00\\
LevenbergMarquardt & 10 & 0.0385 & 3.20 & 4.20 & 4.20\\
MiniBatchGradientDescent & 0 & 0.0618 & 80.00 & 81.00 & 1304.53\\
StochasticGradientDescent & 0 & 0.0394 & 80.00 & 81.00 & 9681.00\\
\bottomrule
\end{tabular}
\end{adjustbox}

\smallnote{SGD и mini-batch дошли до лимита 80 эпох, поэтому не считаются успешными по строгому критерию, хотя SGD на части задач дает близкий loss.}
\end{frame}

\begin{frame}{ЛР5: влияние степени полинома}
\scriptsize
\begin{adjustbox}{max width=\textwidth}
\begin{tabular}{lrrrrr}
\toprule
Данные & degree 1 & degree 2 & degree 3 & degree 4 & degree 5\\
\midrule
near\_linear & 0.03942 & 0.03936 & 0.03928 & 0.03861 & 0.03853\\
nonlinear & 1.17604 & 1.15854 & 0.82320 & 0.80880 & 0.28341\\
\bottomrule
\end{tabular}
\end{adjustbox}

\begin{itemize}
    \item Для near-linear данных высокая степень почти не нужна: первая степень уже описывает основную структуру.
    \item Для nonlinear данных степень 5 резко улучшает loss: модель получает достаточно гибкости для пика и осцилляции.
    \item Это одновременно увеличивает риск переобучения, поэтому дальше исследовалась регуляризация.
\end{itemize}
\end{frame}

\begin{frame}{ЛР5: аппроксимация near-linear данных}
\plotone{PolynomialRegressionObjective__dataset_kind_near_linear_degree_5_n_points_120_noise_variance_0.04_seed_11_x_range_-3.0_3.0_fit.png}
{Истинная зависимость почти линейна. Высокая степень в основном подстраивает малые колебания и шум, поэтому выигрыш по loss небольшой.}
\end{frame}

\begin{frame}{ЛР5: аппроксимация nonlinear данных}
\plotone{PolynomialRegressionObjective__dataset_kind_nonlinear_degree_5_n_points_120_noise_variance_0.09_seed_17_x_range_-2.5_2.5_fit.png}
{Нелинейная зависимость содержит пик и осцилляцию. Полином степени 5 заметно лучше описывает форму данных, чем степени 1--4.}
\end{frame}

\begin{frame}{ЛР5: истории loss}
\plottwo{PolynomialRegressionObjective__dataset_kind_nonlinear_degree_5_n_points_120_noise_variance_0.09_seed_17_x_range_-2.5_2.5_history_loss.png}
        {PolynomialRegressionObjective__dataset_kind_nonlinear_degree_5_n_points_120_noise_variance_0.09_seed_17_x_range_-2.5_2.5_history_loss_by_grad_calls.png}
        {История обучения}
        {Ось по итерациям показывает динамику метода, а ось по вызовам градиента честнее сравнивает вычислительную цену SGD и mini-batch с GN/LM.}
\end{frame}

\begin{frame}{ЛР5: batch size}
\plottwo{lab5_batch_size_comparison_by_epoch.png}
        {lab5_batch_size_comparison_by_grad_calls.png}
        {Размер батча}
        {Малый batch дает частые шумные обновления, большой batch ближе к полному градиенту. Сравнение по эпохам и по $N_g$ разделяет устойчивость и фактическую цену.}
\end{frame}

\begin{frame}{ЛР5: регуляризация nonlinear degree 5}
\scriptsize
\begin{adjustbox}{max width=\textwidth}
\begin{tabular}{lrr}
\toprule
Регуляризация & $f$ & Итерации\\
\midrule
нет & 0.28341 & 4\\
L1, $\lambda=0.001$ & 0.29050 & 4\\
L1, $\lambda=0.01$ & 0.35078 & 5\\
L1, $\lambda=0.1$ & 0.63649 & 53\\
L1, $\lambda=1$ & 0.88710 & 58\\
L2, $\lambda=0.001$ & 0.29969 & 4\\
L2, $\lambda=0.01$ & 0.39649 & 4\\
L2, $\lambda=0.1$ & 0.57235 & 3\\
L2, $\lambda=1$ & 0.69107 & 3\\
Elastic Net, $\lambda=0.001$ & 0.30646 & 4\\
Elastic Net, $\lambda=0.01$ & 0.44208 & 4\\
Elastic Net, $\lambda=0.1$ & 0.65768 & 47\\
Elastic Net, $\lambda=1$ & 0.93095 & 85\\
\bottomrule
\end{tabular}
\end{adjustbox}

\smallnote{Полная оптимизируемая функция растет, потому что штраф входит прямо в $f$. Это не ошибка: регуляризация ограничивает коэффициенты и гибкость модели.}
\end{frame}

\begin{frame}{ЛР5: loss и empirical risk при регуляризации}
\plottwo{lab5_regularization_loss.png}
        {lab5_regularization_empirical_risk.png}
        {Полная цель и риск}
        {Полный loss включает штраф, empirical risk измеряет ошибку аппроксимации. При росте регуляризации модель жертвует точностью на train ради более сдержанных коэффициентов.}
\end{frame}

\begin{frame}{ЛР5: что делает регуляризация с моделью}
\plottwo{lab5_regularization_fit.png}
        {lab5_regularized_coefficients.png}
        {Форма кривой и коэффициенты}
        {L1 сильнее зануляет коэффициенты, L2 уменьшает их плавнее, Elastic Net совмещает оба эффекта. На графике видно, что штраф меняет саму модель, а не только значение функции.}
\end{frame}

\begin{frame}{ЛР5: сравнение методов по loss и времени}
\plottwo{lab5_method_comparison_loss.png}
        {lab5_method_comparison_elapsed_seconds.png}
        {Методы оптимизации}
        {GN и LM используют структуру суммы квадратов и быстро получают малую ошибку. SGD дешев в одном обновлении, но требует много градиентных вычислений и не достиг строгого критерия за 80 эпох.}
\end{frame}

\begin{frame}{ЛР5: выводы для защиты}
\begin{itemize}
    \item Регрессия сведена к минимизации эмпирического риска: loss строится из ошибок предсказания на данных.
    \item Для одномерной линейной модели аналитическое решение самое дешевое, но применимо только в узком случае.
    \item Gauss--Newton и LM эффективны, потому что используют якобиан невязок всей least squares задачи.
    \item LM устойчивее GN за счет демпфирования: плохой шаг увеличивает $\lambda$, хороший шаг уменьшает.
    \item Регуляризация повышает полную цель, но уменьшает гибкость модели и ограничивает коэффициенты.
    \item Строки \texttt{skipped\_batch\_study} в CSV являются техническим фильтром, а не неудачным обучением.
\end{itemize}
\end{frame}

\section{Итог}

\begin{frame}{Связь ЛР4 и ЛР5}
\begin{columns}[T,onlytextwidth]
\column{0.48\textwidth}
\textbf{ЛР4}
\begin{itemize}
    \item классическая безусловная оптимизация;
    \item квадратичные модели;
    \item Гессиан и его аппроксимации;
    \item глобализация через line search и trust-region.
\end{itemize}
\column{0.49\textwidth}
\textbf{ЛР5}
\begin{itemize}
    \item data-driven постановка;
    \item loss строится по выборке;
    \item least squares дает структуру $J^TJ$;
    \item GN и LM используют ту же идею квадратичной модели.
\end{itemize}
\end{columns}

\smallnote{Главная связь: в обеих лабораторных локальная квадратичная модель объясняет направление шага, но в ЛР5 эта модель появляется из невязок данных.}
\end{frame}

\begin{frame}{Короткий чеклист перед защитой}
\small
\begin{itemize}
    \item Почему $H\succ0$ гарантирует направление спуска для Ньютона?
    \item Чем line search отличается от trust-region?
    \item Почему CG имеет конечную сходимость только на SPD-квадратиках?
    \item Что означает условие $s_k^Ty_k>0$ в BFGS/DFP?
    \item Что такое successful run и почему SGD может иметь 0 успешных запусков?
    \item Чем loss отличается от empirical risk и от регуляризованной цели?
    \item Почему GN заменяет точный Гессиан на $J^TJ$?
    \item Почему LM переходит между GN и градиентным спуском?
    \item Что меняют L1, L2 и Elastic Net в коэффициентах модели?
\end{itemize}
\end{frame}

\begin{frame}{Файлы результатов}
\begin{itemize}
    \item ЛР4: \texttt{outputs/lab4/summary.csv}, \texttt{outputs/lab4/plots}.
    \item ЛР5: \texttt{outputs/lab5/tables/summary.csv}, histories, coefficients, datasets.
    \item Графики презентации сохранены в \texttt{presentation/assets/labs\_4\_5/}.
    \item Полные таблицы остаются в CSV; в слайды вынесены агрегаты и ключевые иллюстрации для защиты.
\end{itemize}

\begin{block}{Главный итог}
Методы второго порядка выигрывают, когда Гессиан или структура least squares доступны и надежны; квазиньютоновские и стохастические методы являются компромиссом по памяти, производным и стоимости одной итерации.
\end{block}
\end{frame}

\end{document}
